%COMANDO CHE DETERMINA IL TIPO DI DOCUMENTO CHE SI VUOLE CREARE

\documentclass[a4paper,12pt]{report}

%ELENCO DEI PACCHETTI UTILI PER LA STESURA DEL DOCUMENTO

\usepackage[utf8]{inputenc}
\usepackage[english, italian]{babel}
\usepackage{graphicx}
\usepackage{float}
\usepackage{tabularx}
\usepackage{makecell}
\usepackage{titlesec}
\usepackage{fancyhdr}
\usepackage{lastpage}
\usepackage{xurl}
\usepackage{hyperref}
\usepackage{geometry}
\usepackage[table,dvipsnames]{xcolor}
\setcounter{tocdepth}{5}
\setcounter{secnumdepth}{5}
\usepackage{caption}
\usepackage{etoolbox}% >= v2.1 2011-01-03
\usepackage{tikz}
\usepackage{listings}
\usepackage{scalerel}
\usepackage{longtable}
\usepackage{comment}
\usepackage{amssymb}
\usepackage{listings}

\titleformat{\chapter}[display]
{\normalfont\bfseries}{}{0pt}{\LARGE}

%COMANDO PER LA SPAZIATURA DEI TITOLI DAL BORDO DEL FOGLIO

\titlespacing*{\chapter}{0cm}{0cm}{0.2cm}

%COMANDO PER LA SPAZIATURA DEL TESTO DAI BORDI LATERALI

\geometry{
	left=20mm,
	right=20mm,
}

%COMADNO PER AVERE L'INDICE DEL NOME CHE SI VUOLE

\renewcommand{\contentsname}{Indice}

%COMADNI PER OTTENERE SUBSUBSECTION NUMERATE E PRESENTI NELL'INDICE

\setcounter{tocdepth}{5}
\setcounter{secnumdepth}{5}

%COMANDI PER OTTENERE HEADER E FOOTER

\pagestyle{plain}

\fancypagestyle{plain}{
	\fancyhf{}
	\lhead{\includegraphics[width=3cm]{../../immagini/minilogo.jpg}}
	\chead{}
	\rhead{\fontsize{12}{10}Verbale Esterno 2021-04-15}
	\lfoot{}
	\cfoot{\thepage\ di \pageref{LastPage}}
	\rfoot{}
}

\pagestyle{plain}

%COMANDI PER LINK

\hypersetup{
	colorlinks=true,
	linkcolor=black,
	filecolor=black,
	urlcolor=blue,
	citecolor=black,
}

%COMANDI PER SNIPPET DI CODICE

\BeforeBeginEnvironment{lstlisting}{\begin{mdframed}\vspace{-0.7em}}
	\AfterEndEnvironment{lstlisting}{\vspace{-0.5em}\end{mdframed}}

% needed for \lstcapt
\def\ifempty#1{\def\temparg{#1}\ifx\temparg\empty}

% make new caption command for listings
\usepackage{caption}
\newcommand{\lstcapt}[2][]{%
	\ifempty{#1}%
	\captionof{lstlisting}{#2}%
	\else%
	\captionof{lstlisting}[#1]{#2}%
	\fi%
	\vspace{0.75\baselineskip}%
}

\definecolor{atomlightorange}{rgb}{0.88,0.76,0.55}
\definecolor{atomdarkgrey}{RGB}{59,62,75}

% set listings
\lstset{%
	basicstyle=\footnotesize\ttfamily\color{atomlightorange},
	framesep=20pt,
	belowskip=10pt,
	aboveskip=10pt
}

% add frame environment
\usepackage[%
framemethod=tikz,
skipbelow=8pt,
skipabove=13pt
]{mdframed}
\mdfsetup{%
	leftmargin=0pt,
	rightmargin=0pt,
	backgroundcolor=atomdarkgrey,
	middlelinecolor=atomdarkgrey,
	roundcorner=6
}

%BIBLIOGRAFIA

\makeatletter
\def\thebibliography#1{\chapter*{Bibliografia\@mkboth
		{Bibliografia}{Bibliografia}}\list
	{[\arabic{enumi}]}{\settowidth\labelwidth{[#1]}\leftmargin\labelwidth
		\advance\leftmargin\labelsep
		\usecounter{enumi}}
	\def\newblock{\hskip .11em plus .33em minus .07em}
	\sloppy\clubpenalty4000\widowpenalty4000
	\sfcode`\.=1000\relax}
\makeatother

%DEFINIZIONE COLORI

\definecolor{airforceblue}{rgb}{0.36, 0.54, 0.66}


\begin{document}
	
	\makeatletter
	\begin{titlepage}
		\begin{center}
			\vspace*{-5cm}
			\author{Jawa Druids} 
			\title{Studio di Fattibilità}
			\date{} %LASCIARE QUESTO CAMPO VUOTO, SE LO TOLGO STAMPA LA DATA CORRENTE
			\includegraphics[width=0.5\linewidth]{../immagini/DRUIDSLOGO.jpg}\\[4ex]
			{\huge \bfseries  \@title }\\[2ex] 
			{\LARGE  \@author}\\[50ex]
			\vspace*{-9cm}
			\begin{table}[H]
				\renewcommand{\arraystretch}{1.4}
				\centering
				\begin{tabular}{r | l}
					\textbf{Versione} & 1.0.0 \\%RIGA PER INSERIRE LA VERSIONE ULTIMA DEL DOCUMENTO
					\textbf{Data approvazione} & 02-12-2020\\
					\textbf{Responsabile} & Margherita Mitillo \\
					\textbf{Redattori} & \makecell[tl]{Margherita Mitillo \\ Emma Roveroni} \\
					\textbf{Verificatori} & \makecell[tl]{Igli Mezini \\ Mattia Cocco} \\
					%MAKECELL SERVE PER POI ANDARE A CAPO ALL'INTERNO DELLA CELLA
					\textbf{Stato} & Approvato\\
					\textbf{Lista distribuzione} & \makecell[tl]{JawaDruids \\ Prof. Tullio vardanega \\ Prof. Riccardo Cardin \\ Sync Lab}\\
					\textbf{Uso} & Interno
				\end{tabular}
			\end{table}
			\vspace{0.1cm}
			\hfill \break
			\fontsize{17}{10}\textbf{Sommario} \\
			\vspace{0.1cm}
			Studio di Fattibilità dei capitolati proposti 
		\end{center}
	\end{titlepage}
	\makeatother
	
	\quad
\begin{center}
	\LARGE\textbf{Registro delle modifiche}
\end{center}

\def\tabularxcolumn#1{m{#1}}
{\rowcolors{2}{RawSienna!90!RawSienna!20}{RawSienna!70!RawSienna!40}

\begin{center}
	\renewcommand{\arraystretch}{1.4}
\begin{tabularx}{\textwidth}{|X|c|c|c|c|}
	\hline
	\rowcolor{airforceblue}
	\textbf{Modifica} & \textbf{Autore} & \textbf{Ruolo} & \textbf{Data} & \textbf{Versione}\\
	\hline
	\textit{Descrizione delle modifiche} & Nome Cognome & \textit{Analista} & 23-11-2020 & v0.0.3 \\
	\hline
	\textit{Descrizione delle modifiche} & Nome Cognome & \textit{Analista} & 20-11-2020 & v0.0.2 \\
	\hline
	\textit{ Stesura del documento} & Emma Roveroni & \textit{Analista} & 06-12-2020 & v0.0.1 \\
	\hline
\end{tabularx}
\end{center}
	
	%COMANDO PER LA CREAZIONE DELL'INDICE
	
	\tableofcontents{}
	
	%PER RENDERE PIÙ CHIARA LA STESURA DEI DOCUMENTI È MEGLIO LASCIARE SEPARATI IN FILE DIVERSI OGNI CAPITOLO
	
	% \input{esempio} -- esempio di codice per inserire un nuovo capitolo
\chapter{Introduzione}

\section{Scopo del documento}
Lo scopo di questo documento è la descrizione delle motivazioni che hanno portato il gruppo alla scelta del capitolato C3, ovvero GDP - Gathering Detection Platform. Inoltre sono spiegate anche le motivazioni conseguenti all'esclusione degli altri progetti proposti.

\section{Glossario}
All'interno del documento sono presenti termini che possono risultare ambigui o incongruenti a seconda del contesto in cui si trovano. Per evitare il sorgere di incomprensioni viene fornito un glossario individuabile nel file \textit{Glossario} contenente i suddetti termini con la loro relativa spiegazione.
\\Nella seguente documentazione per favorire chiarezza ed evitare inutili ridondanze tali termini verranno indicati mettendo la lettera "G" come pedice ad ogni prima ricorrenza che si incontra ad inizio di ogni sezione.

\section{Riferimenti}
\subsection{Riferimenti normativi}
\begin{itemize}
	\item \textit{Norme di Progetto}
\end{itemize}
\subsection{Riferimenti informativi}
\begin{itemize}
	\item \textbf{Capitolato d'appalto C1: BlockCOVID, supporto digitale al contrasto della pandemia} \\ 
	\url{https://www.math.unipd.it/~tullio/IS-1/2020/Progetto/C1.pdf}
	\item \textbf{Capitolato d'appalto C2: EmporioLambda, piattaforma di e-commerce in stile Serverless} \\
	\url{https://www.math.unipd.it/~tullio/IS-1/2020/Progetto/C2.pdf}
	\item \textbf{Capitolato d'appalto C3: GDP, Gathering Detection Platform}\\
	\url{https://www.math.unipd.it/~tullio/IS-1/2020/Progetto/C3.pdf}
	\item \textbf{Capitolato d'appalto C4: HD Viz, visualizzazione di dati multidimensionali}\\
	\url{https://www.math.unipd.it/~tullio/IS-1/2020/Progetto/C4.pdf}
	\item \textbf{Capitolato d'appalto C5: PORTACS, piattaforma di controllo mobilità autonoma}\\
	\url{https://www.math.unipd.it/~tullio/IS-1/2020/Progetto/C5.pdf}
	\item \textbf{Capitolato d'appalto C6: RGP, Realtime Gaming Platform}\\
	\url{https://sesaspa-my.sharepoint.com/:b:/g/personal/s\_dindo_vargroup\_it/EThvay0f6KVCoXydYOce2lkBt-MYcnW1yafRXFXVIOIsHg?e=2emZZI}
	\item \textbf{Capitolato d'appalto C7: SSD, soluzioni di sincronizzazione desktop}\\
	\url{https://www.math.unipd.it/~tullio/IS-1/2020/Progetto/C7.pdf}
\end{itemize}
\chapter{Capitolato scelto: C3 - GDP, Gathering Detection Platform}

Il capitolato C3 è stato presentato dall'azienda Sync Lab, Software House  nata dal 2002 che propone nel mercato prodotti software nei settori mobile, videosorveglianza e molti altri. L’obiettivo di Sync Lab è la realizzazione, messa in opera e governance di soluzioni IT ed inoltre è molto sensibile all'innovazione attraverso attività di ricerca  e sviluppo.

\section{Informazioni generali}
\begin{itemize}
	\item \textbf{Nome} - GDP: Gathering Detection Platform
	\item \textbf{Proponente} - Sync Lab S.r.l.
	\item \textbf{Committente} - Prof. Tullio Vardanega e Prof. Riccardo Cardin
\end{itemize}
\section{Descrizione del capitolato}
Il periodo storico che stiamo vivendo ha portato noi cittadini inizialmente ad una quarantena forzata e successivamente è stata concessa la circolazione, ma questo porta alla creazione di situazioni di rischio assembramento e di conseguenza un aumento nel rischio di contagio. L'idea dell'azienda è creare quindi una piattaforma che possa aiutare i cittadini a vivere più serenamente e con sicurezza questa situazione con la quale si possa scoprire quali zone sono più a rischio assembramento rispetto ad altre. \\
Ossia bisogna utilizzare elementi raccolti da sensoristica e da altre sorgenti; in questo modo si può fare una stima ed ottenere indicazioni sui potenziali assembramenti e conseguentemente fornire un supporto per l'ottimizzazione del traffico.
\section{Finalità del progetto}
La soluzione al problema precedentemente descritto deve essere un prototipo software in grado di acquisire e monitorare dati per poi estrapolare informazioni da utilizzare per identificare le possibili zone o i possibili eventi che presentano un alto flusso di utenti. Gli utilizzatori della piattaforma devono poter quindi visualizzare dati in tempo reale tramite heatmap oppure predizioni future di una determinata zona. \\
In particolare gli obiettivi tecnologici che si vogliono raggiungere sono:
\begin{itemize}
	\item realizzazione di un software atto a contare le persone nei mezzi pubblici;
	\item realizzazione di simulazione dei dati per poter monitorare dati storici e previsionali;
	\item capacità di acquisire informazioni a bassa latenza ed in modo continuativo;
	\item elaborazione in tempo reale dei dati;
	\item identificazione di eventi concorrenti;
	\item previsione di variazione del flusso di utenti.
\end{itemize}
Riguardo al lato predittivo si dovrà istruire il software tramite Machine Learnig a riconoscere dati di momenti passati ed elaborare predizioni di future. Inoltre bisogna aggiornare automaticamente le previsioni sulla base dei nuovi dati che vengono osservati.
\section{Tecnologie interessate}
Il proponente è interessato ad esplorare nuove tecnologie quindi predilige non imporre tecnologie specifiche affidandosi alle proposte dei fornitori, ma ha specificato comunque alcune scelte tecnologiche da considerare per lo svolgimento del progetto. Le tecnologie consigliate sono le seguenti:
\begin{itemize}
	\item \textit{Java} e \textit{Angular} per lo sviluppo delle parti di Back-End e Front-End;
	\item il framework \textit{Leaflet} per la gestione delle mappe;
\end{itemize}
\section{Aspetti positivi}
\begin{itemize}
	\item La possibilità di lavorare ad un progetto legato a tematiche contemporanee;
	\item Anche se il progetto nasce da una determinata tematica il gruppo ha concluso che un prodotto software del genere potrebbe essere applicato in diversi campi, quindi vi è una possibilità di riutilizzo dei ragionamenti e dei metodi di sviluppo futuri;
	\item La possibilità di non avere vincoli tecnologici;
	\item Il proponente si è dimostrato aperto al confronto e disponibile a creare un ambiente caratterizzato da una forte collaborazione;
	\item La possibilità di esplorare tecnologie e tematiche non presenti nel percorso di studi universitario;
	\item Interesse da parte di tutti i componenti del gruppo a lavorare con il Machine Learning.
\end{itemize}
\section{Criticità e fattori di rischio}
\begin{itemize}
	\item L'apprendimento delle nuove tecnologie o delle strumentazioni previste potrebbe risultare lento per i membri del gruppo che non le hanno mai utilizzate.
\end{itemize}
\section{Conclusioni}
Il capitolato ha attirato l'attenzione del gruppo fin da subito grazie alla chiarezza e alla linearità della presentazione. Inoltre la possibilità di avere libertà tecnologica e la propensione alla ricerca di tecnologie innovative da parte del proponente potrebbe dare la possibilità di poter imparare ad utilizzare tecnologie non presenti negli insegnamenti del percorso di studi. La presenza del Machine Learning è stato un altro fattore decisivo per la scelta del capitolato, settore interessante che non viene toccato nel percorso di studi. A seguito di queste affermazioni, il gruppo Jawa Druids ha eletto questo capitolato come prima scelta fiduciosi del fatto di riuscire a colmare le lacune e affrontare in sinergia ogni possibile difficoltà che si presenterà nel percorso di creazione del prodotto software richiesto.
\chapter{C1 - BlockCOVID: supporto digitale al contrasto della pandemia}
Il capitolato C1 è stato presentato da Imola Informatica, azienda che si occupa di consulenza IT.
\section{Informazioni generali}
\begin{itemize}
	\item \textbf{Nome} - BlockCOVID: supporto digitale al contrasto della pandemia
	\item \textbf{Proponente} - Imola Informatica SpA
	\item \textbf{Committente} - Prof. Tullio Vardanega e Prof. Riccardo Cardin
\end{itemize}
\section{Descrizione del capitolato}
L'azienda proponente ha deciso di trattare la tematica della pandemia contemporanea in ambito alla sicurezza al lavoro. Infatti ogni azienda dovrebbe assicurare ai propri dipendenti un luogo di lavoro sicuro dal rischio di contagio e, perciò, sanificato correttamente. In particolare devono essere segnalate le postazioni in uso e , successivamente, bisogna comunicare quando vengono liberate in modo che gli addetti possano procedere con una sanificazione corretta, rendendo la postazione pronta per un nuovo utilizzo.\\
Questa procedura, oltre a tutelare il dipendente, tutela anche il datore di lavoro nel caso in cui avvenga un caso di contagio nel luogo di lavoro.
\section{Finalità del progetto}
L'obiettivo è quello di sviluppare un'applicazione in grado di indicare quando una postazione viene occupata da un determinato dipendente. In particolare, tramite l'applicazione, si deve poter sapere se una postazione è libera, occupata oppure prenotata, sapere lo stato di avanzamento della sanificazione e prenotare una postazione. Gli utenti che possono usare questa applicazione sono divisi in tre categorie: amministratore, utente ed addetto alle pulizie. Il primo deve poter gestire le postazioni di lavoro, i dipendenti presenti ed estrapolare un report legato alle postazioni utilizzate da un singolo utente ed uno legato alle sanificazioni effettuate. Il secondo, invece, deve poter prenotare e segnalare l'occupazione della postazione e quando la pulisce con il kit aziendale. Il terzo, infine, deve poter ricevere un elenco delle postazioni che necessitano di sanificazione e spuntare quelle sanificate.
\section{Tecnologie interessate}
Il proponente preferisce non imporre particolari tecnologie da utilizzare per svolgere il progetto in quanto sempre interessato alla ricerca di nuove soluzioni tecnologiche. Di seguito vengono elencate le tecnologie che l'azienda si sente di consigliare al fornitore:
\begin{itemize}
	\item \textit{Java}, \textit{Python} o \textit{nodeJS} per lo sviluppo del Back-End;
	\item \textit{IAAS Kubernetes} oppure un \textit{PAAS} per il rilascio delle componenti.
\end{itemize}
\section{Aspetti positivi}
\begin{itemize}
	\item La possibilità di scegliere autonomamente le tecnologie da utilizzare.
\end{itemize}
\section{Criticità e fattori di rischio}
\begin{itemize}
	\item Il prodotto software di questo capitolato è legato ad una tematica ristretta e che quindi difficilmente potrà essere applicato in campi diversi da quello in cui nasce;
	\item Block Covid è un capitolato che sviluppa un'applicazione con idee già esistenti e nella teoria semplici da capire che non porta il gruppo ad un percorso di autoformazione.
\end{itemize}
\section{Conclusioni}
Questo capitolato, nonostante sia relativo a tematiche contemporanee, non ha suscitato particolare interesse nel gruppo in quanto legato ad una tematica troppo ristretta. Inoltre non si è riusciti ad evidenziare un possibile percorso di autofromazione riguardo a tematiche diverse da quelle del proprio percorso di studi. Per questo motivo, il gruppo ha preferito orientarsi verso un'alternativa più stimolante.
\chapter{C2 - EmporioLambda: piattaforma di e-commerce in stile Serverless}
\section{Informazioni generali}
\begin{itemize}
	\item \textbf{Nome} - EmporioLambda: piattaforma di e-commerce in stile Serverless
	\item \textbf{Proponente} - Red Babel
	\item \textbf{Committente} - Prof. Tullio Vardanega e Prof. Riccardo Cardin
\end{itemize}
\section{Descrizione del capitolato}

\section{Finalità del progetto}
\section{Tecnologie interessate}
\section{Aspetti positivi}
\section{Criticità e fattori di rischio}
\section{Conclusioni}
\chapter{C4 - HD Viz: visualizzazione di dati multidimensionali} \label{CapitolatoC4}
Il capitolato$_{\scaleto{G}{3pt}}$ C4 è stato presentato dalla Zucchetti, società che produce soluzioni software ed hardware e servizi per aziende, assicurazioni e banche.
\section{Informazioni generali} \label{C4InformazioniGenerali}
\begin{itemize}
	\item \textbf{Nome} - HD Viz: visualizzazione di dati multidimensionali;
	\item \textbf{Proponente}$_{\scaleto{G}{3pt}}$ - Zucchetti SpA;
	\item \textbf{Committente}$_{\scaleto{G}{3pt}}$ - Prof. Tullio Vardanega e Prof. Riccardo Cardin.
\end{itemize}
\section{Descrizione del capitolato} \label{C4DescrizioneDelCapitolato}
Con l'evoluzione della tecnologia le applicazioni moderne riescono a memorizzare volumi di dati molto elevati e con molte dimensioni. Anche i programmi tradizionali si sono evoluti in questo campo. L'esposizione grafica di questi dati, però, in certe situazioni, potrebbe risultare inefficiente e causa di errori. \\
Per questo motivo l'azienda richiedere di creare una piattaforma che permetta di visualizzare chiaramente i dati multidimensionali così da poter individuare subito i possibili errori.
\section{Finalità del progetto} \label{C4FinalitàDelProgetto}
In particolare il progetto da realizzare è una piattaforma web dove è possibile visualizzare dati multidimensionali con almeno 15 dimensioni. Questi dati si devono poter fornire sia tramite query che tramite file formato csv. La piattaforma dovrà garantire almeno quattro tipologie di visualizzazioni:
\begin{itemize}
	\item Scatter Plot Matrix, visualizzazione a riquadri posti a matrice;
	\item Force Field, visualizzazione che traduce le distanze tra i punti;
	\item Heat Map$_{\scaleto{G}{3pt}}$, visualizzazione che mostra la distanza tra punti con colori più o meno intensi;
	\item Proiezione Lineare Multi Asse, visualizzazione dei punti in un piano cartesiano.
\end{itemize}
\section{Tecnologie interessate} \label{C4TecnologieInteressate}
Le tecnologie richieste per lo sviluppo di questo progetto sono le seguenti:
\begin{itemize}
	\item HTML, CS e Javascript per lo sviluppo della piattaforma web;
	\item Java$_{\scaleto{G}{3pt}}$ con server Tomcat o in alternativa Javascript con server Node.js per lo sviluppo lato server.
\end{itemize}
\section{Aspetti positivi} \label{C4AspettiPositivi}
\begin{itemize}
	\item La presenza di tecnologie già note potrebbe facilitare la realizzazione del progetto.
\end{itemize}
\section{Criticità e fattori di rischio} \label{C4CriticitàEFattoriDiRischio}
\begin{itemize}
	\item Dalla documentazione sorge che si dovranno utilizzare tecnologie già note ai componenti del gruppo stimolando poco il nostro interesse;
	\item Le tematiche del capitolato$_{\scaleto{G}{3pt}}$ non hanno suscitato interesse particolare in buona parte dei componenti del gruppo.
\end{itemize}
\section{Conclusioni} \label{C4Conclusioni}
Questo capitolato$_{\scaleto{G}{3pt}}$ non ha suscitato interesse nella maggior parte del gruppo, in particolare per via della presenza di tecnologie già note, in quanto già studiate nel percorso di laurea. Per questo motivo il gruppo ha deciso di escluderlo nella scelta del progetto da svolgere.
\chapter{C5 - PORTACS: piattaforma di controllo mobilità autonoma} \label{CapitolatoC5}
Il capitolato$_{\scaleto{G}{3pt}}$ C5 è stato presentato dalla SanMarco Informatica, azienda che offre soluzioni tecnologiche e consulenza digital.
\section{Informazioni generali} \label{C5InformazioniGenerali}
\begin{itemize}
	\item \textbf{Nome} - PORTACS: piattaforma di controllo mobilità autonoma;
	\item \textbf{Proponente}$_{\scaleto{G}{3pt}}$ - SanMarco Informatica SpA;
	\item \textbf{Committente}$_{\scaleto{G}{3pt}}$ - Prof. Tullio Vardanega e Prof. Riccardo Cardin.
\end{itemize}
\section{Descrizione del capitolato} \label{C5DescrizioneDelCapitolato}
L'azienda richiede la realizzazione di un prodotto dove si potranno visualizzare tutte le unità, ossia robot, muletto o automobile, che avranno un punto di partenza in una griglia ed una lista di punti di interesse. Nel lato utente si dovranno vedere quattro frecce direzionali, un pulsante stop/start e un indicatore di velocità. Il sistema dovrà indicare la prossima direzione da prendere in base alla posizione dei punti di interesse e delle altre unità per evitare collisioni.
\section{Finalità del progetto} \label{C5FinalitàDelProgetto}
Il software finito dovrà essere in grado di accettare in input la scacchiera (o mappa), con le percorrenze e i relativi vincoli, e la definizione delle unità, con identificativo di sistema, velocità massima, posizione iniziale e lista di punti di interesse da dover raggiungere già ordinati. \\
Inoltre le unità dovranno inviare costantemente la propria posizione in modo tale che il sistema centrale riesca a pilotare tutte le unità presenti sulla scacchiera.
\section{Tecnologie interessate} \label{C5TecnologieInteressate}
Il proponente$_{\scaleto{G}{3pt}}$ non ha specificato nulla riguardo alle tecnologie da utilizzare per lo svolgimento del progetto.
\section{Aspetti positivi} \label{C5AspettiPositivi}
\begin{itemize}
	\item La possibilità di lavorare per un'azienda grande potrebbe essere vantaggioso per quanto riguarda il curriculum.
\end{itemize}
\section{Criticità e fattori di rischio} \label{C5CriticitàEFattoriDiRischio}
\begin{itemize}
	\item L'assenza di linee guida riguardanti le tecnologie da utilizzare;
	\item La tematica del capitolato$_{\scaleto{G}{3pt}}$ non ha interessato nessuno dei componenti del gruppo;
	\item La documentazione poco chiara e generale.
\end{itemize}
\section{Conclusioni} \label{C5Conclusioni}
Dato che il capitolato$_{\scaleto{G}{3pt}}$ non ha suscitato interesse e la documentazione è risultata poco esplicativa, il gruppo ha deciso di escluderlo nella scelta del progetto da svolgere.
\chapter{C6 - RGP: Realtime Gaming Platform}
\section{Informazioni generali}
\begin{itemize}
	\item \textbf{Nome} - RGP: Realtime Gaming Platform
	\item \textbf{Proponente} - zero12
	\item \textbf{Committente} - Prof. Tullio Vardanega e Prof. Riccardo Cardin
\end{itemize}
\section{Descrizione del capitolato}
Il progetto prevede la realizzazione di un videogioco a scorrimento verticale da utilizzare su dispositivi mobile, con possibilità di giocare sia in modalità single player che multiplayer. La sfida tra più giocatori deve essere la parte centrale dello svolgimento del progetto.
\section{Finalità del progetto}
La parte fondamentale della realizzazione del progetto è lo sviluppo di un'architettura cloud-based  per dare la possibilità di far comunicare diversi dispositivi. La modalità multiplayer deve avere le seguenti caratteristiche:
\begin{itemize}
	\item la sfida deve essere ad eliminazione;
	\item durante la partita deve essere possibile vedere in tempo reale anche gli altri giocatori;
	\item i nemici e il power up devono essere sincronizzati in tempo reale per rendere la sfida equa.
\end{itemize}
Per quanto riguarda la modalità single player invece, questa dovrà essere infinita con livelli di difficoltà sempre crescenti.
\section{Tecnologie interessate}
Il proponente ha esplicitamente indicato le scelte di tecnologiche di seguito elencate.
\begin{itemize}
	\item Swift oppure SwiftUI per lo sviluppo in sistemi operativi IOS;
	\item Kotlin per lo sviluppo in sistemi operativi Android;
	\item NodeJs servizio preferibile per lo sviluppo;
	\item Amazon Web Services per gestire la comunicazione tra diversi device.
\end{itemize}
\section{Aspetti positivi}
\begin{itemize}
	\item Formazione su argomenti al di fuori del percorso di studi.
\end{itemize}
\section{Criticità e fattori di rischio}
\begin{itemize}
	\item Alta complessità di realizzazione del progetto;
	\item Tematica ritenuta interessante solo da una parte del gruppo.
\end{itemize}
\section{Conclusioni}
Nonostante la possibilità di autoformazione riguardo ad argomenti al di fuori del percorso di studi, la tematica non ha suscitato interesse nella maggior parte del gruppo e di conseguenza si è deciso di escluderlo nella scelta del progetto da svolgere.
\chapter{C7 - SSD: soluzioni di sincronizzazione desktop}
\label{CapitolatoC7}
Il capitolato$_{\scaleto{G}{3pt}}$ C7 è stato presentato dall'azienda Zextras, nata come Service Provider Zimbra, cioè un sistema email on-premise
\section{Informazioni generali} \label{C7InformazioniGenerali}
\begin{itemize}
	\item \textbf{Nome} - SSD: soluzioni di sincronizzazione desktop
	\item \textbf{Proponente}$_{\scaleto{G}{3pt}}$ - Zextras
	\item \textbf{Committente}$_{\scaleto{G}{3pt}}$ - Prof. Tullio Vardanega e Prof. Riccardo Cardin
\end{itemize}
\section{Descrizione del capitolato} \label{C7DescrizioneDelCapitolato}
Al giorno d'oggi l'utente necessita di poter accedere ai propri contenuti da diversi dispositivi, sia web che mobile. Il proponente$_{\scaleto{G}{3pt}}$ vorrebbe approfondire la sincronizzazione desktop per poter permettere all'utente il salvataggio in cloud del proprio lavoro.
\section{Finalità del progetto} \label{C7FinalitàDelProgetto}
Il progetto ha come obiettivo finale lo sviluppo di un'interfaccia multipiattaforma che sfrutti un algoritmo efficiente, capace di garantire il salvataggio in cloud del lavoro e la sincronizzazione dei cambiamenti presenti in cloud. Tale interfaccia deve poter essere disponibile per i sistemi operativi più usati (Windows, Mac, Linux), senza domandare all'utente di installare prodotti aggiuntivi per il suo funzionamento. 
\section{Tecnologie interessate} \label{C7TecnologieInteressate}
Il proponente$_{\scaleto{G}{3pt}}$ ha consigliato diverse tecnologie utili per lo svolgimento del progetto e sono le seguenti.
\begin{itemize}
	\item Qt Framework per lo sviluppo dell'interfaccia e del controller dell'architettura;
	\item Python per lo sviluppo del Back-End$_{\scaleto{G}{3pt}}$.
\end{itemize}
\section{Aspetti positivi} \label{C7AspettiPositivi}
\begin{itemize}
	\item Tematica ritenuta interessante per alcuni membri del gruppo.
\end{itemize}
\section{Criticità e fattori di rischio} \label{C7CriticitàEFattoriDiRischio}
\begin{itemize}
	\item La presenza di Qt Framework tra le tecnologie proposte, IDE non apprezzato da alcuni componenti del gruppo;
\end{itemize}
\section{Conclusioni} \label{C7Conclusioni}
Il gruppo ha deciso di escludere dalla scelta del progetto da svolgere tale capitolato$_{\scaleto{G}{3pt}}$ in quanto una parte del gruppo riteneva poco interessante la tematica e per la presenza  del framework$_{\scaleto{G}{3pt}}$ di Qt.

% bibliography, glossary and index would go here.

\end{document}